Our work draws on conversation analysis, multi-agent communication, and message function analysis.

\paragraph{Conversation analysis and grounding.}
Backchannels~\cite{yngve_getting_1970}, continuers~\cite{schegloff_discourse_1982}, and acknowledgment tokens~\cite{jefferson_acknowledgment_1984} are structurally load-bearing in human conversation despite carrying little propositional content~\cite{clark_contributing_1989}; Bavelas~\etal~\cite{bavelas_co_narrators_2000} show that reducing listener backchannels degrades narrative quality, confirming their functional necessity for co-narration. Computational backchannel research spans prediction~\cite{liermann_backchannel_2023, inoue_backchannel_2025} and representation learning~\cite{qian_backchannel_2026}, but focuses on human conversation; whether analogous functional dynamics hold for LLM-generated acknowledgments has not been tested.

\paragraph{Multi-agent communication.}
Classical agent communication languages~\cite{finin_kqml_1994} defined structured message types; learned communication~\cite{sukhbaatar_commnet_2016, foerster_dial_2016, wang_imac_2020} optimized bandwidth end-to-end. Rizvi-Martel~\etal~\cite{rizvi_martel_communication_2026} derive bounds on \emph{when} communication helps; Zheng~\etal~\cite{zheng_thought_comm_2025} study communication \emph{alternatives} (latent thoughts), whereas we study which natural-language messages are safe to \emph{remove}. Chen~\etal~\cite{chen_five_ws_2026} identify ``what to communicate'' as understudied. Multi-agent debate~\cite{du_debate_2024} and LLM frameworks~\cite{li_camel_2023, hong_metagpt_2024, wu_autogen_2024, chen_agentverse_2024, wang_megaagent_2025} demonstrate communication affects performance, though debate may reduce to majority voting~\cite{choi_debate_vote_2025} and intrinsic strength can dominate structure~\cite{wu_can_llm_debate_2025, deliberationbench_2026, la_malfa_miss_mark_2025}; see Yan~\etal~\cite{yan_beyond_self_talk_2025} and Tran~\etal~\cite{tran_collaboration_survey_2025} for surveys. Zhou~\etal~\cite{zhou_why_agents_communicate_2025} and Buscemi~\etal~\cite{buscemi_strategic_comm_2025} study \emph{why} agents communicate but neither performs message-level functional analysis.

\paragraph{Communication efficiency.}
AgentDropout~\cite{wang_agentdropout_2025, wang_agentdropout_v2_2026} prunes entire agents (20--56\% token reduction) but does not distinguish which \emph{messages} are dispensable. S$^2$-MAD~\cite{zeng_s2mad_2025} uses turn-level early stopping; OPTIMA~\cite{chen_optima_2025} optimizes topology and routing; SafeSieve~\cite{safesieve_2025} prunes inter-agent edges via dual-mechanism scoring---none perform message-function analysis. Lin~\etal~\cite{lin_supervisoragent_2026} introduce SupervisorAgent, a lightweight meta-agent that uses LLM-free adaptive filtering to intervene at runtime, reducing token consumption by ${\sim}30\%$; their approach operates at the token/trajectory level without analyzing the communicative \emph{function} of individual messages, whereas our taxonomy-driven pruning targets specific message categories identified through functional analysis. Other approaches span agent pruning~\cite{zhang_agentprune_2025, li_agp_2025, xiao_agentdiet_2026, wang_dagp_2025}, latent-space communication~\cite{zou_latentmas_2025, du_interlat_2025}, KV-cache sharing~\cite{ye_kvcomm_2025}, and holistic optimization~\cite{li_omac_2025, zhang_communication_strategies_2025, jiang_agentcomm_2026} (Appendix Table~\ref{tab:method_comparison}). No prior work has examined whether communication pruning safety is invariant to model scale; we provide the first evidence spanning a $10\times$ parameter range (7B--72B).

\paragraph{Topology and protocol.}
GoAgent~\cite{chen_goagent_2026}, AnyMAC~\cite{wang_anymac_2025}, ProtocolBench~\cite{du_protocolbench_2025}, MultiAgentBench~\cite{zhu_multiagentbench_2025}, LACP~\cite{li_lacp_2025}, GUARDIAN~\cite{zhou_guardian_2025}, and Shen~\etal~\cite{shen_topology_propagation_2025} optimize \emph{who} communicates and \emph{how}; we complement these by analyzing \emph{what} at the message-function level.

\paragraph{Attribution and pragmatics.}
HiveMind~\cite{xia_hivemind_2026}, AgentSHAP~\cite{horovicz_agentshap_2025}, and MAST~\cite{cemri_mast_2025} operate at agent/trace level; our ablation provides message-level granularity. Classical speech act theory~\cite{austin_how_1962, searle_speech_acts_1969} and dialogue act taxonomies~\cite{core_damsl_1997, bunt_iso_2012} inform our categories. Qamar~\etal~\cite{qamar_dialogue_acts_2025} apply dialogue act analysis to LLM conversations; Maitra~\etal~\cite{maitra_dialogue_lens_2025} use dialogue acts as a lens on human--LLM interaction quality. Zhang~\etal~\cite{zhang_communication_strategies_2025} evaluate communication strategies for dual-agent math problem-solving, finding that acknowledgment and hints play key roles. Our work extends this line to inter-agent communication, connecting functional categories to pruning outcomes.
